% Input fragment. Requires the exact semilocal-anchor proposition and the
% finite-fibre compression theorem with one additional variable.

\section{A finite-fiber presentation with 557 retained coordinates}
\label{sec:finite-fibre-application}

We now give a source-specific application of
Theorem~\ref{thm:finite-fibre-one-variable}.  The count is a theorem
about an existential polynomial obtainable by elimination.  It is not
the degree or variable count of an already expanded coefficient table.
Its identification with the ordinary integers remains conditional on
the arithmetic correctness of the source definition, exactly as for
any projection-preserving optimization of that definition.

All geometric fibers in this section are taken after extension to an
algebraic closure of $\mathbb Q$.  This qualification matters:
$\mathbb Q(\sqrt2)\otimes_{\mathbb Q}\overline{\mathbb Q}$ is the
product of two copies of $\overline{\mathbb Q}$, not a field.
Whenever a quadratic-field inverse occurs, we retain the inequation
that its rational norm is nonzero.  Its value is then uniquely determined
on both factors, including on geometric fibers.

\subsection{Point occurrences}

The root module has six nonidentity point occurrences, with labels
$A,S,T,Z$, height $J$, and common comparison $J$.  Each of the 41
comparison calls has five more, with labels $H,X_0,X_1,X_2,X_3$.
Thus the total is
\[
 6+41\cdot5=211.
\]
Replace each such point by the double of a halfpoint $(a,b)$ on
the fixed curve $E_D:b^2=a^3-Da$ over $K=\mathbb Q(\sqrt2)$.
This uses the exact nonidentity $2E_D(K)$ membership statement of the
source.  Retain $a_0,a_1$, where $a=a_0+a_1\sqrt2$, and leave
$b_0,b_1$ as coordinates to be eliminated.  Keep the original
nonzero guards on the doubled point as well as the norm guard on $b$.
All downstream point coordinates are rational functions of these
halfpoint coordinates.

This point substitution can be checked by explicit rational maps,
without a group-membership assertion as a black box.  Write $D=d^2$
for the source curve.  From a halfpoint $(a,b)$ set
\[
 u=\frac{2b}{a^2+D},\qquad
 v=\frac{a^2-2da-D}{a^2+D},\qquad
 w=\frac{a^2+2da-D}{a^2+D},\qquad h=vw.
\]
Then put $f=u^2$, $x=1/u^2$, $y=h/u^3$, and $z=-u/h$.
On the guarded domain these satisfy every old full-chart equation:
\[
 v^2+df=1,\quad w^2+v^2=2,\quad h=vw,\quad
 zh+u=0,\quad xf=1,\quad yz+x=0.
\]
In addition,
\[
 x=\frac{(a^2+D)^2}{4b^2},\qquad
 y=\frac{(a^2+D)(a^4-6Da^2+D^2)}{8b^3},
\]
so $(x,y)=[2](a,b)$ by the duplication formulas.

Conversely, the old full chart forces $u,v,w$ to be nonzero.
With $A=(1+v)(1+w)$ and $S=v+w$, define
\[
 a=\frac{A}{u^2},\qquad b=\frac{SA}{u^3}.
\]
These formulas always give a valid halfpoint.  Indeed
$v^2+w^2=2$ and $du^2=1-v^2$ imply
\[
 A(S^2-A)=-(1-v^2)^2,\qquad
 A^2+(1-v^2)^2=2SA.
\]
The first identity gives $b^2=a^3-d^2a$; the second, together with
\begin{align*}
 A^2-2(1-v^2)A-(1-v^2)^2&=2vSA,\\
 A^2+2(1-v^2)A-(1-v^2)^2&=2wSA,
\end{align*}
shows that the first set of formulas recovers the same $u,v,w$.
None of the new denominators
vanishes.  In fact $1+v=0$ or $1+w=0$ would give $du^2=0$; and
$v+w=0$ would give $v^2=w^2=1$, with the same contradiction.
Thus $A,S,b,a^2+D$ are nonzero.  This proves both directions of the
full-chart substitution directly.  The original guards on the output
exclude any remaining exceptional halfpoints in the forward map.

For every fixed geometric value of $a_0,a_1$, the equation
$b^2=a^3-Da$ has at most four geometric solutions for $(b_0,b_1)$.
Indeed, after the splitting
$K\otimes\overline{\mathbb Q}\simeq
\overline{\mathbb Q}\times\overline{\mathbb Q}$, it becomes two
independent equations $b_+^2=c_+$ and $b_-^2=c_-$, each with at most
two roots.  Restricting to the guarded open subset cannot increase
this number.  The retained point coordinates therefore number
\[
 211\cdot2=422.
\]

\subsection{The pole blocks and their zero fibers}

There is one pole block from the height call and one from each of the
41 comparison calls, hence 42 in total.  The source block has input
$b$ and private variables $t,U,e,p,q,s,h$, with $B=17b$, satisfying
\begin{align*}
 U(1+bt^2)&=b(1-bt^2),\\
 p^2&=B(1-e)(1+3e),\\
 h&=e(B-U^2),\\
 Bq^2&=(s^2-Bh)(s^2+3Bh).
\end{align*}
Retain $t,e,s$ from each block.  The factor $1+bt^2$ never vanishes
on this system, even on geometric fibers.  Otherwise $bt^2=-1$ and
the first equation would give $0=2b$, a contradiction.  Thus $U$ is
uniquely determined; $h$ is then uniquely determined, $p$ has at most
two possibilities, and $q$ has at most two possibilities when $b\ne0$.

The exceptional case $b=0$ must be treated.  The source equations then
force
\[
 U=p=h=s=0,
\]
but they do not constrain the private variable $q$.  Replace this
fiber by $q=0$.  This does not change any rational existential projection:
$q$ occurs nowhere outside this block, and replacing it by zero
preserves all its equations when $b=0$.  The same normalization makes
the geometric fiber finite.  It is a constructible condition, given
by the union of the branch $b\ne0$ and the branch $b=0,q=0$.

Consequently these blocks contribute
\[
 42\cdot3=126
\]
retained coordinates and at most four eliminated choices per block.

\subsection{The smaller anchor}

Let $C\ne0$ be the fixed source anchor constant and let $f_0,f_1$
be the two reciprocal coordinates supplied to the anchor.  Put
\[
 F=f_0^2+f_0f_1-f_1^2.
\]
By Proposition~\ref{prop:explicit-semilocal} and
Lemma~\ref{lem:two-semilocal-inputs}, the exact anchor condition is
equivalent to
\begin{equation}\label{eq:finite-anchor}
 u^2+6v^2-2w^2-3F^2+2C^4=0
\end{equation}
with three rational witnesses.  This follows by applying the
semilocal formula to $F/C^2$ and multiplying by $C^4$; multiplication
of each of its three witnesses by $C^2$ is a bijection of $\mathbb Q$.
Retain $v,w$.  Once the point coordinates and their finite choices
are fixed, $F$ is fixed and \eqref{eq:finite-anchor} is monic quadratic
in $u$.  Thus two retained coordinates suffice for this block, with
at most two eliminated geometric choices.  The old two sets of five
retained anchor coordinates have been replaced by these two.

\subsection{The remaining parameters and algebraic choices}

Retain the seven root coordinates
\[
 X,\quad s,\quad z,\quad\theta,\quad e,\quad u_c,\quad v_c.
\]
Here $X$ is the ordinary input, $s,z$ are two of the three index
proxies, $\theta$ is the common approximation to $\sqrt2$, $e$ is
the common coordinate denominator, and $u_c,v_c$ are the two common
contact parameters.  The third index proxy, denoted $\tau$ to avoid
confusing it with a retained pole variable, satisfies the monic equation
\[
 \tau^2-s(s-2)\tau z+(s-2)^2z^2-1=0.
\]
It therefore contributes at most two geometric choices and no retained
coordinate.  The eight height-class bits satisfy $\epsilon^2=\epsilon$
and contribute only finitely many choices.  Every remaining arithmetic
gate is determined by its inputs; every inverse is unique on its
retained guarded open subset.  Fixed constant constructions contribute
no free coordinates.

One nonvanishing point in this claim deserves explicit attention.
The compressed reciprocal-logarithm equations are
$\Lambda p_0=G_0$, $A_i=C_ip_i$, $A_i\Lambda=G_i$.
Their coefficients $G_i=G_J(Df_i^2)$ do not vanish for $f_i\in K$.
Here
\[
 G_J(W)=1+\sum_{j=1}^J a_jW^j,\quad
 a_1=-2/5,\quad a_j=a_{j-1}\frac{4j-5}{4j+1}\ (j\geq2),
\]
with $J\geq2$ in every application.  The extension of the two-adic
valuation to $K$ is
\[
 v(a+b\sqrt2)=\min\{v_2(a),v_2(b)+\tfrac12\}.
\]
The fixed $D$ has valuation zero and all $a_j$ have valuation one.
If $v(f)\geq0$, the constant term of $G_J(Df^2)$ is uniquely least
in valuation; if $v(f)<0$, then $v(f)\leq-1/2$ and the last
nonconstant term is uniquely least.  Thus $G_J(Df^2)\ne0$.
The displayed equations force $\Lambda,p_i\ne0$ on rational solutions.
Imposing their rational norm guards before eliminating them therefore
loses no rational solutions.  These guards are still retained on
geometric fibers: the valuation argument is not being asserted over
an algebraic closure.

For clarity, the passage from each compressed point chart back to the
full one is also explicit.  Its equations are
\[
 f_i=u_i^2,\quad v_i^2+df_i=1,\quad
 p_i^2=x_i-Df_i,\quad x_if_i=1.
\]
The nonzero $p_i$ just proved and the identity
$f_ip_i^2=v_i^2(1+df_i)$ force $v_i\ne0$.
Set $z_i=-1/p_i$, $y_i=x_ip_i$, $h_i=u_ip_i$, and
$w_i=u_ip_i/v_i$.  Substitution gives all full-chart equations.
Conversely the full chart gives $p_i=y_if_i$.
The ratio variables are then determined by
$\Lambda=G_0/p_0$, $A_i=G_i/\Lambda$, and $C_i=A_i/p_i$.
All other difference and reciprocal-coordinate guards of the original
source are retained.

\begin{proposition}\label{prop:finite-base-557}
After the exact point-membership substitution, the exact anchor
replacement, and the harmless zero-fiber normalization of the private
pole variable $q$, the source has a constructible presentation whose
projection to
\[
 422+126+2+7=557
\]
retained rational coordinates has finite geometric fibers.
Its rational existential projection onto $X$ is unchanged.
\end{proposition}
\begin{proof}
The exact equivalences of the substitutions have been established
above.  In a fiber of the retained coordinates, first choose the
ordinate coordinates of the 211 halfpoints, in finitely many ways.
The doubled points and all their downstream rational functions are
then determined wherever their guards hold.  Next choose the index
proxy $\tau$ and the eight bits.  This determines the rational input
of each pole block.  The normalized pole equations have finitely many
solutions for their remaining private variables by the case analysis
above.  Finally the anchor ordinate $u$ has at most two choices.
Additional equations or inequations only discard choices.
All remaining auxiliary coordinates are unique functions of these
choices.  Thus every geometric fiber is finite, including exceptional
and empty fibers.

For an explicit uniform upper bound on the number of choices, the
argument gives
\[
 4^{211}\cdot4^{42}\cdot2\cdot2\cdot2^8=2^{516}.
\]
This is only a bound on the size of a geometric fiber of this
presentation, not on the degree of the final polynomial after
stratification and branch multiplication.
\end{proof}

\begin{corollary}\label{cor:558-variables}
There is an effectively constructible integer-coefficient polynomial
in at most 558 total variables, one of them the input $X$, whose
projected rational zero set equals that of the source system.
Conditional on that source system defining $\mathbb Z$ in $\mathbb Q$,
the polynomial is a Diophantine definition of $\mathbb Z$ using at
most 557 witnesses.
\end{corollary}
\begin{proof}
Apply Theorem~\ref{thm:finite-fibre-one-variable} with $d=557$ to
Proposition~\ref{prop:finite-base-557}.  Its polynomial uses the 557
retained coordinates and one additional coordinate.  Existentially
quantifying all except $X$ gives the stated projection.
\end{proof}

The corollary is an existence and effective-construction bound, with
explicitly proved projection equivalence.  An actual output polynomial
of this size requires completing the finite-algebra elimination and
the exceptional-stratum computation.  No expanded coefficient table,
degree, coefficient height, or claim of minimum variable count is
asserted here.

